% concatenated from Gurevich.tex

%   See the AIP README file for restrictions and more information.
%
% TeX'ing this file requires that you have AMS-LaTeX 2.0 installed
% as well as the rest of the prerequisites for REVTeX 4.1
% 
% It also requires running BibTeX. The commands are as follows:
%
%  1)  latex  aipsamp
%  2)  bibtex aipsamp
%  3)  latex  aipsamp
%  4)  latex  aipsamp
%

\documentclass[%
 aip,
% jmp,
% bmf,
% sd,
% rsi,
 amsmath,amssymb,
preprint,%
% reprint,%
%author-year,%
%author-numerical,%
% Conference Proceedings
]{revtex4-1}

\usepackage{graphicx}% Include figure files
\usepackage{dcolumn}% Align table columns on decimal point
\usepackage{bm}% bold math
%\usepackage[mathlines]{lineno}% Enable numbering of text and display math
%\linenumbers\relax % Commence numbering lines

\usepackage[utf8]{inputenc}
\usepackage[T1]{fontenc}
%\usepackage[tbtags]{amsmath}
%\usepackage{amsfonts,amssymb,mathrsfs,amscd}
\usepackage{mathptmx}
\usepackage{etoolbox}

%-------------------------------------------------

\newtheorem{theorem}{Theorem}
\newtheorem{lemma}{Lemma}[section]
\newtheorem{propos}{Proposition}
\newtheorem{cor}{Corollary}

%-------------------------------------------------

\newtheorem{definition}{Definition}
%\newtheorem{proof}{Proof}\def\theproof{}
\newtheorem{remark}{Remark}
\newenvironment{proof}{{\bf Proof:}}
{\hfill $\diamond$\medskip}


%% Apr 2021: AIP requests that the corresponding 
%% email to be moved after the affiliations
\makeatletter
\def\@email#1#2{%
 \endgroup
 \patchcmd{\titleblock@produce}
  {\frontmatter@RRAPformat}
  {\frontmatter@RRAPformat{\produce@RRAP{*#1\href{mailto:#2}{#2}}}\frontmatter@RRAPformat}
  {}{}
}%
\makeatother
\begin{document}

%\preprint{AIP/123-QED}

\title[Three heteroclinic orbits induce a countable family of equivalence classes of regular flows]{Three heteroclinic orbits induce a countable family of equivalence classes of regular flows}
% Force line breaks with \\
\author{E. Gurevich}
 \affiliation{NRU HSE, Nizhny Novgorod}%Lines break automatically or can be forced with \\
\email{egurevich@hse.ru}
\date{\today}% It is always \today, today,
             %  but any date may be explicitly specified

\begin{abstract}
We solve the problem of topological classification for smooth structurally stable flows on closed four-dimensional manifolds, the non-wandering set of which  contains exactly two saddle equilibria, and the wandering set contains isolated trajectories connecting these saddle equilibria (heteroclinic curves). In particular, we show that for a flow of the class under consideration on $\mathbb{CP}^2$, the number of heteroclinic curves is a complete topological invariant, while on  the sphere $\mathbb S^4$,  there exists a countably many  equivalence classes with  an arbitrary odd number $\gamma\geq 3$ of heteroclinic curves. These results contrast with a  three-dimensional case, where under similar conditions there exists only finite set of equivalence classes for each number of heteroclinic curves. 
\end{abstract}

\maketitle

\begin{quotation}
The problem of topological classification of gradient-like flows, that are structurally stable flows with finite nonwandering set,  has a long and rich history, originating in the classical works of A. A. Andronov, L. S. Pontryagin, E. A. Leontovich, and A. G. Mayer~[\onlinecite{AP37,LM37,LM55}]. For gradient-like flows on closed manifolds of dimensions 2 and 3, this problem was completely solved in~[\onlinecite{Pe, Um90}]. For the  dimension of the  ambient  manifolds equal to  four  and higher, the most complete results have been obtained only for gradient-like flows without  heteroclinic  trajectories, connecting saddle equilibria  (see~[\onlinecite{Pil, MeZh13, MeZh16, GrGu22, GrGu23, GuSa24, GuSa25, GuSa26}]). The goal of this paper is to study the effects of  heteroclinic trajectories on the asymptotic behavior of gradient-like flows trajectories on four-dimensional manifolds. 
\end{quotation}



\section{Introduction\label{intro}}

Let  $M^n$ denote a smooth closed manifold of dimension~$n\geq 1$. Recall that an equilibrium state $p$ of a smooth flow $f^t$ on $M^n$ is called {\it hyperbolic} if the Jacobi matrix $$J_p=(\frac{\partial F}{\partial x})|_{_p}$$ of the   velocity field $F=\frac{\partial f^t(x)}{\partial t}|_{_{t=0}}$ at the point $p$ has no eigenvalues with zero real part. The sets $$W^s_p=\{x\in M^n: \lim\limits_{t\to +\infty} f^{t}(x)= p\},$$ $$W^u_p=\{x\in M^n: \lim\limits_{t\to +\infty} f^{-t}(x)= p\}$$ are called {\it stable and unstable invariant manifolds of the equilibrium state $p$}, respectively. The number $\mu_p$ equal to the dimension $\dim W^u_{p}$ of the unstable manifold $W^u_p$ is called {\it the  Morse index of $p$}. If  $i_p=0(n)$ then $p$ is called {\it a sink (a source, respectively)}, and  if $i_p\in (0,n)$, $p$ is called {\it a saddle}. 

A smooth flow $f^t$ on $M^n$ is called {\it gradient-like} if its nonwandering set is finite and consists of hyperbolic equilibria, and the invariant manifolds of the equilibria intersect each other  transversely. If $p,q$ are saddle equilibria such that $W^s_p\cup W^u_q\neq \emptyset$, we will call the intersection $W^s_p\cup W^u_q$ and any trajectory in it  {\it heteroclinic}.  

Let $G_{\mu,\nu,k}$ denote the class of gradient-like flows on closed four-dimensional manifolds whose set of saddle equilibria consists of exactly two points with Morse indices $\mu, \nu\in \{1,2,3\}$, $\mu\leq \nu$. The number $k$ denotes the total number of equilibria of such flows. The following theorem describes all possible combinations of the numbers $\mu, \nu, k$ and the topology of the manifold $M^4$ admitting flows from the corresponding class.

\begin{theorem} \label{top} Let $f^t\in G_{\mu,\nu,k}$. Then exactly one of the following implications holds.

\begin{enumerate}\itemsep=-2pt
\item $k=6, \mu, \nu\in \{1,3\}$ or $k=4$, $\mu\in \{1,2\}, \nu=\mu+1$, and $M^4$ is homeomorphic to the sphere $\mathbb S^4$.

\item $k=5$, $\mu\in \{1,2\}, \nu=\mu+1$, and $M^4$ is homeomorphic to the complex projective plane $\mathbb{CP}^2$.

\item $k=4, \mu=1, \nu=3$, and $M^4$ is homeomorphic to either the direct product $\mathbb S^3\times \mathbb S^1$, or to the unique non-orientable $\mathbb S^3$-bundle over $\mathbb S^1$.
\item $k=4, \mu=\nu=2$, and $M^4$ is homeomorphic to one of the following manifolds: the connected sum $\mathbb{CP}^2\sharp \mathbb{CP}^2$, the direct product $\mathbb S^2\times \mathbb S^2$, or a non-trivial $\mathbb S^2$-bundle over $\mathbb S^2$.

\end{enumerate}
\end{theorem}

For $\mu=\nu$, any flow from the class under consideration has no heteroclinic intersections. Each of the classes described in Theorem~\ref{top}, except $G_{1,2,4}\cup G_{2,3,4}$, contains flows without heteroclinic intersections, the topological classification of which follows from~[\onlinecite{Pil,GrGu23, GrGu22, GuSa24}]. Below, in Theorems~\ref{clas+},~\ref{clas}, and \ref{realiz}, we obtain necessary and sufficient conditions for the topological equivalence of flows from the classes $G_{1,2,k}$ and describe an algorithm for constructing a representative of each  class of topological  equivalence. The topological classification of flows from classes $G_{2,3,k}$ reduces to theorems~\ref{clas+}, \ref{clas} by time reversing. In other cases, the problem of topological classification can be solved using the same methods and is of no particular interest.

\begin{theorem}\label{clas+} If the wandering set of a flow $f^t\in G_{1,2,5}$ contains a heteroclinic intersection, then it consists of an even number of heteroclinic trajectories. Flows $f^t, {f'}^t\in G_{1,2,5}$ are topologically equivalent if and only if they either have no heteroclinic intersections or have the same number of heteroclinic trajectories. For each number $\gamma\geq 0$, there exists a flow $f^t\in G_{1,2,5}$ that has exactly $2\gamma$ heteroclinic trajectories.
\end{theorem}

We show that for the class $G_{1,2,4}$ a similar result is true only if the number of heteroclinic curves  is equal to   one. Let us describe a topological invariant that determines a topological equivalence class of  a flow  $f^t\in G_{1,2,4}$. Due to Theorem~\ref{top}, the ambient manifold of $f^t$ is a topological sphere $S^4$.  Let $\sigma_1, \sigma_2$ be saddle equilibria of $f^t$ with Morse indices equal to 1 and 2, respectively, and $$\Sigma\subset S^4$$ be a smooth closed three-dimensional submanifold of the sphere $S^4$ that separates the sets $A={\rm cl}\,W^u_{\sigma_1}$ and $R={\rm cl}\,W^s_{\sigma_2}$,  and transversely intersects each trajectory belonging to  $S^4\setminus (A\cup R)$. We will call $\Sigma$ {\it a characteristic cross-section} of the flow $f^t$. Set $$\Lambda_s=\Sigma\cap W^s_{\sigma_1}, \lambda_u=\Sigma\cap W^u_{\sigma_2}.$$ 

We refer to  a  triple $\{\Sigma, \Lambda_s, \lambda_u\}$  as  {\it a scheme of a flow $f^t$}. We show in Lemma~\ref{seq1} that for any flow $f^t\in G_{1,2,4}$ the manifold $\Sigma$ is diffeomorphic to the direct product $\mathbb S^2\times \mathbb S^1$, $\Lambda_s$ is a homotopically  nontrivial two-dimensional sphere in $\Sigma$, and $\lambda_u$ is a trivial knot, i.e., a smoothly embedded simple closed curve such that $g(\lambda_u)=\{x\}\times S^1$ for some diffeomorphism $g: \Sigma\to S^2\times S^1$. 

\begin{theorem}\label{clas} Any flow $f^t\in G_{1,2,4}$ has a nonempty set of heteroclinic curves consisting of an odd number of curves. Flows $f^t, {f'}^t\in G_{1,2,4}$, having a unique heteroclinic trajectory, are topologically equivalent. Flows $f^t, {f'}^t\in G_{1,2,4}$ that have more than one heteroclinic trajectory are topologically equivalent if and only if there exists a homeomorphism $h: \Sigma\to \Sigma'$ such that $h(\Lambda_s)=\Lambda_s'$, $h(\lambda_u)=\lambda_u'$.
\end{theorem}


\begin{theorem}\label{realiz} Let $\Lambda, \lambda \subset \mathbb S^2\times \mathbb S^1$ be  smoothly embedded homotopically nontrivial sphere and the trivial knot, respectively. Then there exists a flow $f^t\in G_{1,2,4}$ whose characteristic  cross-section $\Sigma$ is homeomorphic to $\mathbb S^2\times \mathbb S^1$ via a homeomorphism $h:\Sigma\to \mathbb S^2\times \mathbb S^1$ such that $h(\Lambda_s)=\Lambda$ and $h(\lambda_u)=\lambda$.

In particular, for every integer $\gamma \geq 1$, the class $G_{1,2,4}$ contains a countable set of topologically nonequivalent flows with exactly $(2\gamma+1)$ heteroclinic trajectories.
\end{theorem}
%%%%%%%%%%%%%%%%%%%%%%%%%%%%%%%%%%%%%%

\section{Topological Aspects of the Dynamics of Flows from the Class $G_{\mu, \nu, k}$}\label{top_s}

In this section, we prove Theorem~\ref{top} and Lemma~\ref{seq1}, which describe the  topology of the  embedding of  closures of invariant manifolds of saddle equilibria.

\subsection{Energy Function}

Recall that a $C^2$-smooth function $\varphi: M^n\to \mathbb R$ is called a {\it Morse function} if all its critical points are non-degenerate. The {\it Morse index} of a non-degenerate critical point $p\in M^n$ is the number of negative eigenvalues of the Hessian matrix evaluated at this point.
In~[\onlinecite{Sm61,Me68}], it was proven that for any gradient-like flow  there exists {\it an energy function},  a Morse function that is  strictly decreasing along  trajectories other than equilibrium states and has a critical point at each equilibrium state. This implies the following assertion, which refines the properties of the energy function for flows in the class under consideration.


\begin{propos}\label{SM}
For any flow $f^t\in G_{\mu, \nu, k}$, there exists an energy function, i.e., a Morse function $\varphi: M^4\to [0,4]$ with the following properties:
\begin{enumerate}
\item $\varphi$ is strictly decreasing along flow trajectories other than equilibrium states;
\item the trajectories of the flow $f^t$ transversely intersect the set $\varphi^{-1}(c)$ for any non-critical value $c\in (0.4)$ of the function $\varphi$;
\item the set of critical points of the function $\varphi$ coincides with the set $\Omega_{f^t}$;
\item for any equilibrium $p\in \Omega_{f^t}$ such that ${\rm dim}\,W^u_p=\mu_p\in \{0,\dots, 4\}$,
there exists a neighborhood $U_p$ equipped with local coordinates (Morse coordinates) $y_1,\dots, y_4$ such that the sets $W^u_p\cap U_p, W^s_p\cap U_p$ are described by the equations
$y_{\mu_p+1}=\dots=y_4=0$, $y_1=\dots=y_{\mu_p}=0$, respectively, and the function $\varphi$ has the form 
$$\varphi(y_1,\dots, y_4)=\mu_p+\sum\limits_{i=1}^4\alpha_iy_i^2,$$
where $\alpha_i\in \{+1, -1\}$ and the number of negative $\alpha_i$ is equal to $\mu_p$.
\end{enumerate}
\end{propos}

\subsection{Proof of the Theorem~\ref{top}}

We provide the proof of the theorem for the case $\mu=1,\nu=2$, which can be easily generalized to the remaining cases except for $\mu=\nu=2$ and the cases where $k=4,\mu=1,\nu=3$. The theorem for these last cases was proven in~[\onlinecite{GuSa24,GuSa25}]. 

Let $f^t\in G_{1, 2, k}$, $\sigma_1, \sigma_2$ be the saddle equilibria of the flow $f^t$ of Morse indices 1, 2, respectively, $M^4$ be the supporting manifold, and $\varphi: M^4\to [0,4]$ be the energy function of the flow $f^t$ satisfying the conclusion of Proposition~\ref{SM}. It was shown in~[\onlinecite[Corollary 10.1.8, p. 225]{Pa88}] that if $M^n$ is not simply connected, then every Morse function on $M^n$ has a critical point of index 1. Since $\psi=4-\varphi$ has no such critical points, $M^4$ is simply connected and hence orientable.

Since the function $\varphi$ decreases along the trajectories of the flow $f^t$, the manifold $M^4$ can be represented as unions of stable manifolds of all equilibria of the flow $f^t$. It follows from~[\onlinecite[Theorem 8,9]{HW}] that if the complement of $M^4$ to some subset $X$ is disconnected, then the dimension of $X$ is three. Since the flow $f^t$ has a unique equilibrium state $\sigma_1$, whose stable manifold is three-dimensional, the set $M^4\setminus W^s_{\sigma_1}$ has at most two connected components. Consequently, only one of two possibilities is realized: 1) the flow $f^t$ has a unique sink $\omega$; 2) the flow $f^t$ has two sinks $\omega_1, \omega_2$. Moving on to the function $\psi=4-\varphi$, which is the energy function for the flow $f^{-t}$, we similarly find that the set of equilibrium states of the flow $f^t$ contains exactly one source $\alpha$. Then $k=4$ in case 1) and $k=5$ in case 2).

Let $\beta_i$ denote the $i$-th Betti number of the manifold $M^4$. Since $M^4$ is simply connected, $\beta_0=1, \beta_1=0$. It follows from Poincare duality  that $\beta_4=\beta_0=1, \beta_3=\beta_1=0$. By the Morse equality (see~[\onlinecite[\S 5]{Morse}]) we have $$k_0-k_1+k_2-k_3+k_4=\beta_0-\beta_1+\beta_2-\beta_3+\beta_4.$$ Then in case 1) $\beta_2=0$ and $M^4$ is a simply connected homology sphere. It follows from ~[\onlinecite[Theorem~1.6]{Fr}] that $M^4$ is homeomorphic to the sphere $S^4$. In case 2) we obtain $\beta_2=1$. It follows from~[\onlinecite[Theorem~1.6]{Fr}] that there exists a unique (up to homeomorphism) simply connected smooth manifold with such a Betti number, and this manifold is the complex projective plane $\mathbb{CP}^2$. Theorem~\ref{top} is proven.

\begin{remark}\label{nw} In the proof of Theorem~\ref{top}, it was shown that the nonwandering set of the flow $f^t\in G_{1,2,k}$, in addition to two saddles $\sigma_1, \sigma_2$, contains exactly one source and one (two) sinks in the case $k=4$ ($k=5$, respectively).

\end{remark}

\subsection{Handle decomposition of the ambient manifold}\label{dec}

Everywhere below, unless otherwise stated, an {\it $n$-ball or $n$-ball} is a manifold $B^n$ homeomorphic to the unit ball $$\mathbb B^n=\{x\in \mathbb R^n|\, |x|\leq 1\},$$ an open $n$-ball and an {\it $(n-1)$-sphere} $S^{n-1}$ are manifolds homeomorphic to the interior ${\rm int}\, \mathbb B^n$ and the boundary $\partial \mathbb B^n$ of  $\mathbb B^n$, respectively.

Recall that a {\it handle of index $i$} or just  {$i$-handle} is the direct product $$H^n_i=\mathbb {B}^i\times \mathbb {B}^{n-i}.$$ Disks $\mathbb B^i\times \{O\}$ and $\{O\}\times \mathbb B^{n-i}$ are called the {\it foot and secant  disks} of the $i$-handle, respectively. The boundary of the handle is naturally represented as the union of two subsets $$P_i=\partial \mathbb B^i\times \mathbb B^{n-i}, \, Q_i=\mathbb B^i\times \partial\mathbb B^{n-i}$$ with a common boundary. The set $P$ is called a {\it foot of the handle}. Spheres $\partial{\mathbb B}^i\times \{O\}$ and $\{O\}\times \partial{\mathbb B}^{n-i}$ are called the {\it foot  and secant spheres}, respectively.

Let $X$ be an $n$-dimensional compact manifold with boundary, $\psi: P_i \to \partial X$ be a smooth embedding, and $$Y=X\cup_\psi H^n_i$$ be the manifold obtained by identifying points $x\in P_i$ and $\psi(x)$. We say that $Y$ is obtained by {\it gluing the handle $H^n_i$} to the manifold $X$.
Representing a manifold as a union of handles is called its {\it handle decomposition}.

The Morse function $\varphi: M^n\to M^n$ defines a handle decomposition  of $M^n$ into handles (see, for instance, [\onlinecite[Theorems 3.2, 3.4]{M}]) such that each $i$-handle contains a unique critical point of index $i$. This, together with Proposition~\ref{SM} and Remark~\ref{nw}, implies the following assertion, which describes the handle decomposition of the ambient  manifold $M^4$ of a flow $f^t\in G_{1,2,k}$,  determined by its energy function $\varphi$.

For $c\in [0,4]$, we set $M_c=\varphi^{-1}([0,c])$, $\Sigma_c=\partial M_c$. Note that if $c$ is a regular value of $\varphi$, then $\Sigma_{c}$ is a smooth closed three-dimensional submanifold of $M^4$ that transversally intersects the trajectories of the flow $f^t$. We choose numbers $c_1, c_2, c_3$ such that
\begin{equation}\label{const}
0<c_1<1<c_2<2<c_3<4.
\end{equation}

\begin{propos}\label{hd}
\begin{enumerate}
\item The sets $M_{c_1}$ and $M^4\setminus {\rm int}\,  M_{c_4}$ are the union of 0-handles and a 4-handle, respectively.
\item For each $i\in \{1,2\}$, the set $M_{c_{i+1}}$ is obtained from $M_{c_{i}}$ by gluing the $i$-handle.

\item The manifold
$\Sigma_{c_2}$ is diffeomorphic to $\mathbb S^2\times \mathbb S^1$ for $k=4$ and to the sphere $\mathbb S^3$ for $k=5$.
\end{enumerate}
\end{propos}
\begin{proof}
From Proposition~\ref{SM} and Theorem~\ref{top}, taking into account Remark~\ref{nw}, it follows that the function $\varphi$ has a unique maximum point and one (two) minimum points in the case $k=4$ ($k=5$, respectively). Therefore, the manifold $M^4$ is obtained from $M_{c_3}$ by gluing a single handle $M^4\setminus {\rm int}\, M_{c_3}$ of index 4. It follows that $\Sigma_{c_3}$ is a $3$-sphere and the manifold $M_{c_3}$ is connected. On the other hand, $M_{c_3}$ is obtained from $M_{c_2}$ by gluing a handle of index 2. Since gluing handles of index 2 does not change the number of connected components of the manifold, the manifold $M_{c_2}$ is also connected. In turn, in the case $k=4$, the manifold $M_{c_2}$ is the $0$-handle  $M_{c_1}$ with a 1-handle glued on it. By Theorem~\ref{top}, the manifold $M^4$ is homeomorphic to $\mathbb {CP}^2$, and hence is orientable. Therefore, $M_{c_2}$ is homeomorphic to the direct product $\mathbb S^1\times \mathbb B^3$, and its boundary $\Sigma_{c_2}$ is homeomorphic to $\mathbb S^1\times \mathbb S^2$. In the case $k=5$, the manifold $M_{c_2}$ is a connected union of two 0-handles and one 1-handle, and hence is homeomorphic to the ball $\mathbb B^4$. Thus, in this case $\Sigma_{c_2}$ is homeomorphic to the sphere $\mathbb S^3$. Homeomorphic manifolds of dimension 3 are also diffeomorphic, from which we obtain the required statements.
\end{proof}

\subsection{Properties of the  scheme of the flow  $f^t\in G_{1,2,k}$}

In this section, we refine the definition of the scheme of the  flow in the introduction for flows in the class $G_{1,2,4}$, extend it to the case of flows in the class $G_{1,2,5}$, and clarify some properties of the schemes. 


\begin{definition}\label{seq}
We say that a closed three-dimensional topological submanifold $\Sigma\subset M^4$ is a cross-section for a flow $f^t$ defined on $M^4$ if there exists a topological embedding $e: \Sigma\times [-1,1]\to M^4$ such that $e(\Sigma\times \{0\})=\Sigma$
and for any point $x\in \Sigma$, the segment $e([-1,1]\times \{x\})$ belongs to the trajectory of the point $x$. If each trajectory of the flow $f^t$ intersects $\Sigma$ at a unique point, then the cross-section  $\Sigma$ is called a global cross-section.
\end{definition}

Let $f^t\in G_{1,2,k}$, $k\in \{4,5\}$, $\sigma_1, \sigma_2$ be the saddle equilibria of the flow $f^t$ of Morse indices 1, 2, respectively, $M^4$ be the ambient manifold, and $\varphi: M^4\to [0,4]$ be the energy function of the flow $f^t$ satisfying the conclusion of Proposition~\ref{SM}. For an arbitrary $c_2\in (1,2)$, we set $$ \Sigma=\varphi^{-1}(c_2), \Lambda_s=\Sigma\cap W^s_{\sigma_1}, \lambda_u=\Sigma\cap W^u_{\sigma_2}.$$

\begin{definition}\label{sh-exdef} The set $S_{f^t}=\{\Sigma, \Lambda_s, \lambda_u\}$ %, \tilde \lambda_u\}$
is called the scheme of the flow  $f^t\in G_{1,2,k}$.
\end{definition} 

\begin{definition}\label{sh-eq} Schemes $S_{f^t}=\{\Sigma, \Lambda_s, \lambda_u\}$, $S_{{f'}^t}=\{\Sigma', \Lambda'_s, \lambda'_u\}$ of  flows $f^t, {f'}^t\in G_{1,2,k}$ are called equivalent if there exists  a homeomorphism $h: \Sigma\to \Sigma'$ such that $h(\Lambda_s)=\Lambda_s'$, $h(\lambda_u)=\lambda_u'$. 
\end{definition}

Let us remark that  $\Sigma$ is a global cross-section for the restriction of $f^t$ to the set $M^4\setminus ({\rm cl}\, W^u_{\sigma_1}\cup {\rm cl}\, W^s_{\sigma_2})$.

\begin{propos}\label{indep} The definition of the scheme does not depend on the choice of the energy function $\varphi$ and the secant $\Sigma$. Namely, if $\widehat \Sigma$ is a topological submanifold that is a global cross-section for the restriction of the flow $f^t$ to a set $M^4\setminus ({\rm cl}\, W^u_{\sigma_1}\cup {\rm cl}\, W^s_{\sigma_2})$ different from $\Sigma$, and $\widehat \Lambda_s=\widehat \Sigma\cap W^s_{\sigma_1}, \widehat \lambda_u=\widehat \Sigma\cap W^u_{\sigma_2}$, then the scheme $\{\widehat{\Sigma}, \widehat \Lambda_s, \widehat \lambda_u\}$ is equivalent to the scheme $\{\Sigma, \Lambda_s, \lambda_u\}$.
\end{propos}
\begin{proof} The homeomorphism $h$, which realizes the equivalence of schemes, associates with each point $x\in \widehat \Sigma$ a point $h(x)$ lying in the intersection of the orbit $\mathcal O_x$ of $x$ and the secant $\Sigma$.
\end{proof}

By Proposition~\ref{hd}, the manifold $\Sigma$ is diffeomorphic either to the sphere $\mathbb S^3$ or to the direct product $\mathbb S^3\times \mathbb S^1$.
Recall that a simple closed curve $\lambda\subset S^3$ is called a \textit{trivial knot} if there exists a locally flat embedding $e: \mathbb B^2\to S^3$ such that $e(\partial \mathbb B^2)=\lambda$. We refer to  a simple closed curve $\lambda \subset S^2\times S^1$ as a trivial knot in $S^2\times S^1$ if there  exists a homeomorphism $h: S^2\times S^1\to \mathbb S^2\times \mathbb S^1$ such that $h(\lambda)=\{x\}\times \mathbb S^1$, $x\in \mathbb S^2$. By the well-known "lightbulb theorem" (see, e.g., [\onlinecite[Ex.3 \S F Ch.9]{Ro}]), a smooth knot $\lambda\subset S^2\times S^1$ is trivial if and only if there exists a sphere $S^2\times \{x\}$, $x\in S^1$, such that the intersection $\lambda\cap (S^2\times \{x\})$ consists of a single point. 

\begin{lemma}\label{seq1} Let $f^t\in G_{1,2,k}$. Then $\lambda_u$ is a trivial knot smoothly embedded in $\Sigma$, and $\Lambda_s$ is a 2-sphere smoothly embedded in $\Sigma$. If $k=4$, then $\Lambda_s$ does not divide $\Sigma$.
\end{lemma}
\begin{proof} By definition, $\Lambda_s=\Sigma\cap W^s_{\sigma_1}$ is a closed two-dimensional manifold that is global cros-section for the restriction of $f^t$ to $W^s_{\sigma_1}\setminus \sigma_1$. Then  $\Lambda_s$ is a strong deformation retract of the set $W^s_{\sigma_1}\setminus \sigma_1$ and is homotopically  equivalent to it. Since the equilibrium state of $\sigma_1$ is hyperbolic, $W^s_{\sigma_1}\setminus \sigma_1$ is homeomorphic to $\mathbb R^3\setminus \{O\}$, and, consequently,  is homotopically equivalent to the sphere $\mathbb S^2$. Therefore, $\Lambda_s$ is diffeomorphic to $\mathbb S^2$.

In what follows, we use the notation of Section~\ref{dec}. In the proof of Proposition~\ref{hd}, we established that the manifold $M_{c_2}$ is obtained from $M_{c_1}$ by gluing a single 1-handle $H^4_1$. The set $W^s_{\sigma_1}\cap M_{c_2}$ is a 3-ball that divides the handle $H^4_1$. In the case $k=4$, the manifold $M_{c_2}$ is the result of gluing the handle $H^4_1$ to the ball $H^4_0=M_{c_1}$. Consequently, $W^s_{\sigma_1}\cap M_{c_2}$ does not divide $M_{c_2}$, and $\Lambda_s=W^s_{\sigma_1}\cap \partial M_{c_2}$ does not divide $\Sigma=\partial M_{c_2}$.

Now let us  prove that the knot $\lambda_u$ is trivial. Recall that $M_{c_2}=\varphi^{-1}([0,c_2])$, $c_2\in (1,2)$. Set $N_{c_2}=M^4\setminus {\rm int}\,M_{c_2}$.
The function $\psi=4-\varphi$ is the energy function for the flow $f^{-t}$ and $N_{c_2}=\psi^{-1}([0,4-c_2])$. From the properties of the Morse function, Proposition~\ref{SM} and Theorem~\ref{top}, it follows that $N_{c_2}$ is obtained by gluing a 2-handle to the ball $N_{3.5}=\psi^{-1}([0,0.5])$. Therefore, $\partial N_{c_2}$ is obtained from the sphere $\partial N_{3.5}$ by an  integer Dehn  surgery along the knot $\lambda=\partial N_{3.5}\cap W^s_{\sigma_2}$. But $\partial N_{c_2}=\Sigma$, and $\Sigma$, as proved above, is diffeomorphic either to $\mathbb S^2\times \mathbb S^1$ or to $\mathbb S^3$. It is proven in~[\onlinecite[Theorems~3,2]{GoLu}] that nontrivial Dehn surgery along nontrivial knot never yields the   sphere or  $\mathbb S^2\times S^1$. Then, the knot $\lambda$ is trivial in $\partial N_{3.5}$. The manifolds $\partial N_{3.5}\setminus \lambda$ and $\Sigma\setminus \lambda_u$ are diffeomorphic by means of a  diffeomorphism that assigns to a point $x\in \partial N_{3.5}\setminus \lambda$ the  point $\mathcal O_x\cap (\Sigma\setminus \lambda_u)$. By [\onlinecite[Theorem~1]{GoLu}), knots on  the sphere are determined by their complements.  For the case $k=5$ this immediately implies that  the knot $\lambda_u$ is trivial.

Consider the case $k=4$. In this case, the knot $\lambda$ belongs to the sphere $\partial N_{3.5}$, and the knot $\lambda_u$ belongs to the manifold $\Sigma$, which is diffeomorphic to $\mathbb S^2\times \mathbb S^1$. Let $P\subset \Sigma$ be a compact tubular neighborhood of the knot $\lambda_u$. By what was proved above, the manifold $Q=\Sigma\setminus {\rm int}\,P$ is diffeomorphic to the complement in the sphere $\partial N_{3.5}$ to an open tubular neighborhood of the trivial knot $\lambda$. Hence, $Q$ is diffeomorphic to a solid torus. There is only one way to obtain $\mathbb S^2\times \mathbb S^1$ by gluing two solid tori, in which case the gluing homeomorphism maps the meridian of the first solid torus to the meridian of the second. Therefore, the meridians of the solid tori $P, Q$ coincide. Let $x\in \mathbb S^2$, $D\in \mathbb S^2$ be a disk containing the point $x$ inside, $P_0=D\times \mathbb S^1$, $Q_0=(\mathbb S^2\setminus {\rm int}\, D)\times \mathbb S^1$, and $h_0: P\to P_0$ be a homeomorphism such that $h_0(\lambda_u)=x\times \mathbb S^1$. Since $h_0$ is a homeomorphism of solid tori, it maps the meridian $P$ to the meridian $P_0$. It follows that the homeomorphism $h_0|_{\partial P}=h_0|_{\partial Q}$ extends to a homeomorphism $h_1: Q\to Q_0$ such that $h_1|_{\partial Q}=h_0|_{\partial Q}$ (see, for example, [~\onlinecite[Ch.2, \S 5, Ex.5]{Ro}]). Then the homeomorphism $h: \Sigma_{f^t}\to \mathbb S^2\times \mathbb S^1$, defined by the formula
$$h(x)=\begin{cases}h_0(x), x\in P\cr h_1(x), x\in Q, \end{cases},$$ maps the knot $\lambda_u$ to the knot $x\times \mathbb S^1$. Therefore, the knot $\lambda_u$ is trivial. The lemma is proven.
\end{proof}

Note that in case $k=5$,  Jordan-Brown Theorem and Generalized Schoenflies Theorem imply that the two-dimensional sphere $\Lambda_s$ divides the sphere $\Sigma$ into two connected components, the closure of each of which is a three-dimensional ball.

Let $\lambda, \Lambda\subset \Sigma$ be a smoothly embedded 1-and 2-spheres that  intersect each other transversely. We fix the orientation of all three manifolds. If, at a point $x\in \lambda\cap \Lambda$, the orienting tangent frames to $\lambda, \Lambda$ generate a tangent space to $\Sigma$ whose orientation coincides with the chosen orientation of $\Sigma$, then we assign $+1$ to the point $x$. Otherwise, we assign $-1$ to the point $x$. The sum of all such numbers for all points in $\lambda\cap \Lambda$ is called the {\it intersection index} of the pair $\lambda, \Lambda$. The intersection index is a homological invariant. Therefore, it follows directly from Lemma~\ref{seq1} that the intersection index of the knot $\lambda_u$ and the sphere $\Lambda_s$ is equal to one for $k=4$ and  to zero for $k=5$. The trajectory of any point $x\in \lambda_u\cap \Lambda_s$ is the heteroclinic trajectory lying in the intersection $W^u_{\sigma_2}\cap W^s_{\sigma_1}$. Thus, the following statement is proven.

\begin{cor}\label{nmb_of_het} If $f^t\in G_{1,2,k}$ and $k=4(5)$, then the flow $f^t$ has an odd $($even$)$ number of heteroclinic trajectories.
\end{cor}

\begin{cor}\label{inters1} The schemes of flows $f^t, {f'}^t\in G_{1,2,4}$  each of which  have exactly one heteroclinic trajectory  are equivalent.
\end{cor}
\begin{proof} Since $\Lambda_s$ and $\lambda_u$ intersect  each other transversely at a unique point, there exist compact tubular neighborhoods $N_{_{\Lambda_s}}, N_{\lambda_u}\subset \Sigma$ of the sphere $\Lambda_s$ and the knot $\lambda_u$, respectively, such that the intersection $N_{_{\Lambda_s}}\cap N_{\lambda_u}$ is diffeomorphic to $B^2\times [0,1]$, and the fibers of the two-dimensional and one-dimensional foliations defining the structure of the direct product lie on the fibers of the foliations transversal to $\lambda_u, \Lambda_s$, determined by the structure of the tubular neighborhoods $N_{_{\Lambda_s}}, N_{\lambda_u}$ (See, for example, [\onlinecite[Ch.4, Section 6, Exercise 2]{Hi}]).

Since the characteristic cross-section $\Sigma$ is diffeomorphic to $\mathbb S^2\times \mathbb S^1$, there exists a covering $p:\mathbb R^3\setminus \{O\}\to \Sigma$. Since the sphere $\Lambda_s$ does not divide $\Sigma$, its complete preimage $p^{-1}(\Lambda_s)$ is a countable set of pairwise disjoint smoothly embedded 2-spheres, each of which bounds a ball in $\mathbb R^3$ whose interior contains the point $O$. Let $A\subset \mathbb R^3$ be the compact subset bounded by the pair of spheres $S_0, S_{1}\subset p^{-1}(\Lambda_s)$ such that ${\rm int}\, A\cap p^{-1}(\Lambda_s)=\emptyset$. Set $$\widehat N_{\Lambda_s}=p^{-1}(N_{\Lambda_s})\cap A,\  \widehat N_{\lambda_u}=p^{-1}(N_{\Lambda_s})\cap A,$$ $$\widehat D=A \setminus {\rm int}(\widehat N_{\Lambda_s}\cup \widehat N_{\lambda_u}).$$ The boundary of $\widehat D$ is a piecewise smooth sphere $S^2$. The Generalized Schoenflies Theorem implies that $\widehat D$ is homeomorphic to the ball $\mathbb B^3$. We set $D=p(\widehat D)$. Since the restriction of $p$ to ${\rm int A}$ is a homeomorphism, $D$ is homeomorphic to the ball.

We denote by $N_{_{\Lambda'_s}}, N_{\lambda'_u}, D'\subset \Sigma'$ the sets corresponding to the flow ${f'}^t$ similar  to the sets $N_{_{\Lambda_s}}, N_{\lambda_u}, D\subset \Sigma$. Let $$\xi: N_{_{\Lambda_s}}\cup  N_{\lambda_u}\to N_{_{\Lambda'_s}}\cup  N_{\lambda'_u}$$ be a diffeomorphism such that
$\xi(\Lambda_s)=\Lambda'_s, \xi(\lambda_u)=\lambda'_u$. By Alexander Theorem, any homeomorphism from the boundary of a ball can be extended to the ball, so the homeomorphism $\xi$ can be extended to a homeomorphism $\Xi:D\to D'$ such that $\Xi|_{\partial D}=\xi|_{\partial D}$. Therefore, there is a homeomorphism $h: \Sigma\to \Sigma'$ such that $h(\Lambda_s)=\Lambda_s'$, $h(\lambda_u)=\lambda_u'$, which means that the flow schemes $f^t, {f'}^t$ are equivalent.
\end{proof}

One more corollary of Lemma~\ref{seq1}  can be proved similarly to~[\onlinecite[Theorem~4]{MeZh13}]. 

\begin{cor}\label{flat} The closure ${\rm cl}\, W^s_{\sigma_2}$ of the set $W^s_{\sigma_2}$ is a locally flat 2-sphere in $M^4$.
\end{cor}

We note that in~[\onlinecite{MeZh13, GuSa24}], an algorithm for constructing a series of polar flows with two saddle equilibria, the closures of whose invariant manifolds are wild two-dimensional spheres. 



\section{Necessary and Sufficient Conditions for Topological Equivalence of Flows from $G_{1,2,k}$}

In this section, we prove that the flow scheme $f^t\in G_{1,2,k}$ introduced in Definition~\ref{sh-exdef} is a complete topological invariant of it, that is, the following theorem holds.

\begin{theorem}\label{cond} Flows $f^t, {f'}^t\in G_{1,2,k}, k\in \{4,5\}$, are topologically equivalent if and only if their schemes are equivalent. % there exists a homeomorphism $h: \Sigma\to \Sigma'$ such that $h(\Lambda_s)=\Lambda_s'$, $h(\lambda_u)=\lambda_u'$.
\end{theorem}

Theorem~\ref{clas} is an immediate consequence of Theorem~\ref{cond} and Corollaries~\ref{nmb_of_het},~\ref{inters1}. Theorem~\ref{clas+} follows from Theorem~\ref{cond}, Lemma~\ref{cp2}, proved below, and Lemma~\ref{rea}, the proof of which (independent of the proof of Theorem~\ref{cond}) is presented in the next section.

The necessity of the conditions of Theorem~\ref{cond} follows directly from the definition of topological equivalence and Proposition~\ref{indep}. The proof of sufficiency is presented below in the sequence of Propositions~\ref{Step1}-\ref{Step4}.

\subsection{Two auxiliary topological lemmas}

\begin{lemma}\label{cp2} Let $\Lambda, \Lambda'\subset S^3$ be two-dimensional spheres smoothly embedded in $S^3$, and let $\lambda, \lambda'\subset S^3$ be smooth trivial knots such that the intersection of the sphere $\Lambda (\Lambda')$ and the knot $\lambda (\lambda')$ is transversal and consists of $b (b')$ points. If $b=b'$, then there exists a homeomorphism $h: S^3\to S^3$ such that $h(\Lambda)=\Lambda', h(\lambda)=\lambda'$.
\end{lemma}
\begin{proof} Since the intersection $\lambda\cap \Lambda$ is transversal, $b$ is even. Let $b=0$. Denote by $B_{\lambda} (B_{\lambda'})$ the connected component of $S^3\setminus \Lambda (S^3\setminus \Lambda')$ that contains the knot $\lambda (\lambda')$. Since the knots $\lambda, \lambda'$ are trivial, there exists a homeomorphism $h: S^3\to S^3$ such that $h(\lambda)=\lambda'$. From~[\onlinecite[Theorems 3.1, 3.2 of Chapter 8]{Hi}] it follows that for any two balls with locally flat boundaries in $S^3$ there exists an ambient isotopy with compact support that maps one ball to the other. Therefore, we can assume that $h(S^3\setminus B_\lambda)=S^3\setminus B_{\lambda'}$ (a detailed proof of this fact is presented, for example, in~[\onlinecite[Proposition~4]{GuSa26}]). Thus, for the case $b=0$, the lemma is proved. We prove the lemma for the case $b=2k>0$.

It follows from  Generalized Schoenflies Theorem that the closures $B_+, B_-$ of the connected components of the set $\Sigma\setminus \Lambda$ are three-dimensional balls.
We orient the knot $\lambda$ arbitrarily and number the points $p_1,\dots, p_b$ belonging to the intersection $\lambda\cap \Lambda$ in the order they appear when traversing the knot $\lambda$ in the chosen direction. We denote by $N_\lambda$ the solid torus that is a tubular neighborhood of the knot $\lambda$ in $S^3$. Since the intersection $N_\lambda\cap \Lambda$ is transversal, without loss of generality we can assume that $N_\lambda\cap \Lambda$ consists of a finite set of pairwise disjoint two-dimensional balls $D_1,\dots, D_b$ such that $p_i\subset {\rm int}D_i$ for each $i\in \{1,\dots, b\}$. Set $H_+=B_+\cup N_\lambda, H_-=B_-\cup N_\lambda$. Each of the compact manifolds $H_{\pm}$ is a ball with $k$ handles of index 1, i.e., a handlebody of genus $k$.

Let us show that the triviality of the knot $\lambda$ implies that the closure  $R_+(R_-)$ of the complement to $H_{+} (H_-)$ is also a handlebody of genus $k$. Let $\widetilde \lambda\subset \partial N_\lambda$ be a knot whose linking number   with  $\lambda$ is zero and which  intersect $\Lambda$ transversely at $b$ points $\tilde p_1,\dots, \tilde p_b$. We suppose that  the indices of points $\{\tilde p_i\}$ are chosen in the same way as the indices of the points $\{p_i\}$. For each $i\in \{1,3,\dots, b-1\}$ we  connect the points $\tilde p_i, \tilde p_{i+1}$ by an arc $\gamma_i\subset \Lambda$ such that $\gamma_i\cap \gamma_j=\emptyset$ for $i\neq j$. In the same way, we connect  the points $\tilde p_i, \tilde p_{i+1}$ by an arc $\tilde l_i\subset \tilde \lambda\cap B_-$ and set  $\mu_i=\tilde l_i\cup \gamma_i$.  Then  the union  $\{\mu_i\}$ is a trivial link (a set of trivial knots, each of which lies in a ball that has no common points with other knots in the link).  Since the link is trivial, there exist locally flatly embedded in $B_-$ pairwise disjoint two-dimensional balls $P_1, P_3,\dots, P_{b-1}$ such that $\partial P_i=\mu_i$, $i\in \{1,3,\dots, b-1\}$. Let $Q_1,Q_3,\dots, Q_{b-1}\subset B_-$ be pairwise disjoint cylindrical neighborhoods of these balls, that is, the images of  embeddings $e_i$ of the 2-handle $H^3_2=\mathbb B^2\times \mathbb B^1$ into $B_-$ such that $e_i(\mathbb B^2\times \{0\})=P_i$  (see Fig.~\ref{hb1}). Then the sets $X=H_+\cup (\bigcup \sum \limits_{j=1}^kQ_{2j-1})$, $Y=S^3\setminus {\rm int}\, X$ are locally flat three-dimensional balls in $S^3$, and $R_+=Y\cup (\bigcup \sum \limits_{j=1}^kQ_{2j-1})$ is a ball with $k$ glued 1-handles. It can be shown similarly that $R_-$ is also a handlebody of genus $k$.

\begin{figure}
\centering{\includegraphics[width=0.25\textwidth]{hb.eps}}
\caption{Sphere $\Lambda$, trivial knot $\lambda$, its neighborhood $N_\lambda$, and disk $P_1$}
\label{hb1}
\end{figure}

Objects similar to those introduced above, but corresponding to the pair $\lambda', \Lambda'$ will be denoted by the same symbols with an added prime. To prove the lemma, it suffices to construct a homeomorphism $\Phi_+:\partial H_+\to H'_+$ such that $\Phi_+(\Lambda)=\Lambda', \Phi_+(\lambda)=\lambda'$, and extendable to $H_+\cup R_+$. Recall that the meridians of a handlebody $H$ of genus $k$ are defined as a set of pairwise disjoint curves $c_1,\dots, c_k\subset \partial H$, each of which bounds an open disk in ${\rm int}\, H$, and $\partial H\setminus \bigcup\limits_{i=1}^k c_i$ is homeomorphic to the sphere $S^2$ with $k$ disks cut out. A homeomorphism defined on the boundary of a handlebody extends to a homeomorphism of the hanldebody $H$ if and only if it preserves the  meridians. Note that in the case of handlebodies $H_+$ and $R_+$, the meridians of the first body are the boundaries of the disks $D_1, D_3, \dots D_{b-1}$, and the meridians of the second are the boundaries of the disks $P_1, P_{3},\dots, P_{b-1}$.



Let $\psi:N_\lambda\to  N_{\lambda'}$ be a diffeomorphism such that $\psi(\lambda)=\lambda', \psi(\tilde \lambda)=\tilde \lambda'$, $\psi(N_{\lambda}\cap B_+)\subset B'_+$, and $\psi(D_i)=D'_i$ for each $i\in \{1,\dots, b\}$. From~[\onlinecite[Theorem 3.2 Chapter 8]{Hi}] it follows that there exists a homeomorphism $\Psi: \Lambda\to \Lambda'$ such that $\Psi|_{D_i}=\psi|_{D_i}$ for any $i\in \{1,\dots, b\}$.

The homeomorphisms $\Psi, \psi|_{\partial N_\lambda}$ determine homeomorphisms $\Phi_+: \partial H_{+}\to \partial H'_+, \Phi_-: \partial H_{-}\to \partial H'_-$ as follows:
\begin{equation}
\Phi_{\pm}(x)=\begin{cases} \Psi(x), x\in \Lambda;\cr
\psi(x), x\in \partial N_\lambda.
\end{cases}
\end{equation}

From the definition of the homeomorphisms $\psi, \Phi_{+}$ it follows that the set of curves $\{\Phi_{+}(\partial Q_i)\}$ forms a trivial link. Then there exists a set of pairwise disjoint locally flatly embedded two-dimensional balls $\widetilde P_1,\widetilde P_3, \dots, \widetilde P_{b-1}$ such that $\partial \widetilde P_i=\{\Phi_+(\partial Q_i)\}$ for each $i\in\{1,3,\dots, b-1\}$. Hence,  these curves can be viewed as the meridians of the handlebody $R'_+$, and the homeomorphism $\Phi_+$ extends to a homeomorphism $\Theta+:R_+\to R_+$. Similarly, it can be proved that the homeomorphism $\Phi_-$ extends to a homeomorphism $\Theta_-:R_-\to R_-$. Then the desired homeomorphism $h:S^3\to S^3$ is defined by the relations

\begin{equation}
h(x)=\begin{cases} \psi(x), x\in N_\lambda;\cr
\Theta_+(x), x\in R_+=B_-\setminus {\rm int}\,N_\lambda;\cr
\Theta_-(x), x\in R_-=B_+\setminus {\rm int}\, N_\lambda.
\end{cases}
\end{equation}
\end{proof}

\begin{lemma}\label{dehn}
Let $\mathbb B^2_r=\{x\in \mathbb R^2|\ |x|=r\}$, $0<r_1<r<1$; $x\in \partial \mathbb B^2_r$ --- some point and $g: \mathbb B^2_r\times [0,1]\to \mathbb B^2_r\times [0,1]$ --- a homeomorphism such that $g(\{O\}\times [0,1])=\{O\}\times [0,1]$. Then there exists a homeomorphism $h:\mathbb B^2\times [0,1]\to \mathbb B^2\times [0,1]$ with the following properties:
\begin{enumerate}
\item $h|_{\mathbb B^2_{r_1}\times [0,1]}={\rm id}$;
\item $h(\mathbb B^2_r\times [0,1])=\mathbb B^2_r\times [0,1]$;
\item $h(g(x\times [0,1]))=x\times [0,1]$;
\item $h|_{\partial \mathbb B^2\times [0,1]}={\rm id}$.
\end{enumerate}

\end{lemma}
\begin{proof} Set $l=x\times [0,1], l'=g(l)$. Let $e:[-1,1]\times [0,1]\to \partial \mathbb B^2_r\times [0,1]$ be a topological embedding such that $e(\{0\}\times [0,1])=l$. Set $N_l=e([-1,1]\times [0,1])$ and denote by $e', N_{l'}$ a similar embedding and its image, which is a compact neighborhood of the arc $l'$ in $\partial \mathbb B^2_r\times [0,1]$.
The set $$C_l={\rm cl} (\partial \mathbb B_r\times [0,1]\setminus N_l)$$ is homeomorphic to the disk $\mathbb B^2$. By Alexander Theorem, the restriction of the homeomorphism $e'e^{-1}: N_l\to N_{l'}$ to $\partial N_l$ extends to a homeomorphism $C_l\to C_{l'}$. Therefore, there exists an orientation-preserving homeomorphism
$$\psi: \partial \mathbb B_r\times [0,1]\to \partial \mathbb B^2_r\times [0,1]$$
such that $\psi(l)=l'$. 

Set $$P_+=(\mathbb B^2\setminus {\rm int}\, \mathbb B^2_{r}) \times [0,1],\  P_-=(\mathbb B^2_{r}\setminus {\rm int}\, \mathbb B^2_{r_1}) \times [0,1].$$

 $P_+, P_-$ are homeomorphic to the solid torus   $\mathbb B^2\times \mathbb S^1$. We extend $\psi$ to a homeomorphism $\Psi_\pm:\partial P_\pm\to \partial P_\pm$ that will be identity on $\partial B_r\times [0,1], \partial B_{r_1}\times [0,1]$ and could be extended to a homeomorphism $\Phi_\pm: P_\pm\to P_\pm$. 

Let $\psi_{i}=\psi|_{\partial \mathbb B^2_r\times \{i\}}, i\in \{0,1\}$, and denote by $\psi^t_i: \partial \mathbb B^2_r\times \{i\}\to \partial \mathbb B^2_r\times \{i\}$ an isotopy connecting $\psi^1_i=\psi_i$ with the identity map $\psi^0_i$.  Set $$\tau_+(|x|)=\frac{1-|x|}{1-r},$$ and  define the   homeomorphisms $\Psi_+: \partial P_+\to \partial P_+$ by the relations

$$
\Psi_+(x)=\begin{cases}\psi(x),\ x\in \partial \mathbb B^2_r\times [0,1];\cr
x,\ x\in \partial \mathbb B^2\times [0,1];\cr
\frac{|x|}{r} \psi^{\tau_+(|x|)}_i(\frac{rx}{|x|}),\ x\in \mathbb B^2\setminus \mathbb B^2_r\times \{i\}, i\in \{0,1\}.
\end{cases}
$$

Let  $\rho_{x}$  be a radial  segment of   $\mathbb B^2$ joining the point $x\in \partial \mathbb B^2_r$ with a point  $x'\in \partial \mathbb B^2$. Then the union 

$$\mu=l\cup \{\rho_x\times \{0,1\}\}\cup \{x'\times [0,1]\}$$

is the meridian of solid torus $P_+$. A curve $\Psi_+(\mu)\subset \partial P_+$ is a closed curve that intersects a canonical longitude $\partial \mathbb B^2\times \{0\}$ of the torus $\partial P_+$ at exactly one point. Then there  are two possibilities. 

1)   $\Psi_+(\mu)$ is a meridian of $P_+$, in this case there exists a homeomorphism $\Phi_+: P_+\to P'_+$ such that $\Phi_+|_{\partial P_+}=\Psi_{+}|_{\partial P_+}$.

2) There is a Dehn twist $\tau: \partial P_+\to \partial P_+$ with a support $(\mathbb B^2\setminus {\rm int}\, \mathbb B^2_r)\times \{0\}$ such that
$\tau(\Psi_+(\mu))$ is  a meridian of $P_+$.  In this case the composition $\tau \Psi_+$ sends $l$ to $l'$ and  extends to the desired homeomorphism $\Phi_+: P_+\to P'_+$.

 Similarly, one may construct a homeomorphism $\Phi_-: P_-\to P_-$ such that $\Phi_-|_{\partial P_+\cap \partial P_-}=\Phi_+|_{\partial P_+\cap \partial P_-}=\psi_{\partial P_+\cap \partial P_-}$. Then the  desired homeomorphism $h$ is defined by the formula

\begin{equation}
h(x)=\begin{cases}\Phi_+(x), x\in (\mathbb B^2\setminus {\rm int}\, \mathbb B^2_{r}) \times [0,1];\cr
\Phi_-(x), x\in (\mathbb B^2_r\setminus {\rm int}\, \mathbb B^2_{r_1}) \times [0,1];\cr
x,\ x\in \mathbb B_{r_1}\times [0,1].\end{cases}
\end{equation}

\end{proof}

\subsection{Refining the properties of the homeomorphism $h: \Sigma\to \Sigma'$ satisfying the definition of equivalence of schemes}

From the transversality of the intersection $\Lambda_s\cap \lambda_u$, it follows that there exist tubular neighborhoods $N_{_{\Lambda_s}}, N_{\lambda_u}\subset \Sigma$ of the sphere $\Lambda_s$ and the knot $\lambda_u$, respectively, such that the intersection $N_{_{\Lambda_s}}\cap N_{\lambda_u}$ consists of a finite set of connected components each of which is diffeomorphic to $B^2\times [0,1]$, and the fibers of the two-dimensional and one-dimensional foliations on this component, defining the direct product structure, lie on the fibers of foliations transverse to $\lambda_u, \Lambda_s$ that defined by the structure of tubular neighborhoods $N_{_{\Lambda_s}}, N_{\lambda_u}$ (see, for example,~[\onlinecite[Ch.4, Section 6, Exercise~2]{Hi}]). Since the secant $\Sigma$ is an orientable manifold, the neighborhood $N_{\lambda_u}$ of the knot ${\lambda_u}$ is a solid torus. Let $\tilde \lambda_u\subset \partial N_{\lambda_u}$ denote the longitude  of the solid torus $N_{\lambda_u}$ whose intersection with each connected component of the set $N_{_{\Lambda_s}}\cap N_{\lambda_u}$ consists of a single one-dimensional arc, which is a fiber of the one-dimensional foliation.

Let $N_{_{\Lambda'_s}}, N_{\lambda'_u}$, $\widetilde \lambda'_u$ be similar manifolds for a flow ${f'}^t$ whose scheme is equivalent to the scheme of the flow $f^t$.

\begin{propos}\label{Step0} Let the schemes of the flows $f^t, {f'}^t\in G_{1,2,k}$ be equivalent. Then there is a homeomorphism $\eta: \Sigma\to \Sigma'$ such that:
\begin{enumerate}
\item $\eta(\Lambda_s)=\Lambda'_s$, $\eta(\lambda_u)=\lambda'_u$,
\item $\eta(N_{_{\Lambda_s}})=N_{_{\Lambda'_s}}$, $\eta(N_{\lambda_u})=N_{\lambda'_u}$;
\item if $k=5$, then $\eta(\widetilde \lambda_u)=\widetilde \lambda'_u$.
\end{enumerate}
\end{propos}
\begin{proof} We prove the proposition for $k=4$ (the arguments for $k=5$ are similar). In this case, by Proposition~\ref{hd} and Lemma~\ref{seq1}, the manifold $\Sigma$ is diffeomorphic to $S^2\times S^1$, and the sphere $\Lambda_s$ does not divide $\Sigma$.

By the condition, the schemes of the flows  $f^t, {f'}^t\in G_{1,2,4}$ are equivalent,  that  means the  existing of  a homeomorphism $h: \Sigma\to \Sigma'$ such that $h(\Lambda_s)=\Lambda_s'$, $h(\lambda_u)=\lambda_u'$ (see definition~\ref{sh-eq}). We show that the existence of $h$ implies the existence of the homeomorphism $\eta$.

Set $C=\Sigma\setminus {\rm int}\, N_{\Lambda_s}, C'=\Sigma'\setminus {\rm int}\, N'_{\Lambda_s}$. The proof consists of two steps. In the first step, we show that there exists an orientation  preserving homeomorphism $\eta_{_{C,C'}}: C\to C'$ such that $\eta_{_{C,C'}}(\lambda_u\cap C)=\lambda'_u\cap C$,
$\eta_{_{C,C'}}(N_{\lambda_u}\cap C)=N_{\lambda'_u}\cap C'$, and
$\eta_{_{C,C'}}(\widetilde \lambda_u\cap C)=\widetilde \lambda'_u\cap C$. In the second step, we show that the homeomorphism $\eta_{_{C,C'}}$ can be extended to the set $N_{\Lambda_s}$  to  the desired homeomorphism $\eta$.

{\it Step 1. Constructing the homeomorphism $\eta_{_{C,C'}}: C\to C'$}

Consider the covering $p:\mathbb R^3\setminus \{O\}\to \Sigma$.
Since the sphere $\Lambda_s$ does not divide $\Sigma$, its complete preimage $p^{-1}(\Lambda_s)$ is a countable set of pairwise disjoint smooth  two-dimensional spheres  each of which bounds a ball in $\mathbb R^3$ whose interior contains the point $O$. Let $A\subset \mathbb R^3$ be the compact subset bounded by the pair of spheres $S_0, S_{1}\subset p^{-1}(\Lambda_s)$ such that ${\rm int}\, A\cap p^{-1}(\Lambda_s)=\emptyset$. Set $\widehat N_{\Lambda_s}=p^{-1}(N_{\Lambda_s})\cap A$. Note that one of the connected components of $p^{-1}(C)$ lies in ${\rm int}\, A$ and the restriction of $p$ to this component is a diffeomorphism onto $C$. We will identify this connected component with $C$ denoting it by the same symbol, and keep the notation $\lambda_u\cap C$, etc. for sets $p^{-1}(\lambda_u)\cap C$, etc. Note that Annulus Theorem implies that the sets $A, C$ are diffeomorphic to the annulus $\mathbb S^2\times [0,1]$, as is each connected component of the set $\widehat N_{\Lambda_s}$. Denote by $N_{c}$ the collar of the boundary $\partial C$ in $C$, that is, the image of a smooth embedding $e: \mathbb S^2\times [0,1]\times \mathbb S^0\to C$ such that $ e(\mathbb S^2\times \{0\}\times \mathbb S^0)=\partial C$. Set $N=\widehat N_{\Lambda_s}\cup N_c$ and denote by $N_{-1}, N_{1}$ the connected components of the set $N$. We will define an arbitrary point $p\in N_{\pm 1}$ by two coordinates $x\in \partial C, y\in [0,1]$  and assume that $(x,0)\subset \partial A, \widehat N_{\Lambda_s}\cap N_{\pm 1}=\{(x,y)|\, y\in [0,1/2]\}$.

We define a diffeomorphism $\psi: A\to C$ by setting
\begin{equation}
\psi(p)=\begin{cases}p, p\in A\setminus {\rm int}\, N;\cr
(x, 0.5y+0.5), p=(x,y)\in N_{\pm 1}
\end{cases}
\end{equation}

Let $p',\psi', A', C'$ denote  mappings and sets similar to those introduced above, but corresponding to the flow ${f'}^t$. The homeomorphism $h: \Sigma\to \Sigma'$ defines a homeomorphism $$\tilde h: A\to A'$$ such that $p'\tilde h|_A=h p|_{A}$. We define the homeomorphism $g_{_{C,C'}}: C\to C'$ setting  $$g_{_{C,C'}}=\psi\tilde h\psi^{-1}.$$ By construction, $g_{_{C,C'}}(\lambda_u\cap C)=\lambda'_u\cap C'$, $g_{_{C,C'}}(N_c)=\widetilde N'_c$.

Due to  results of Bing and Moise,see~[\onlinecite[Theorem 8]{Bi}], a compact locally flat subset of a closed 3-manifold is a subpolyhedron, so the compact arcs $\lambda'_u\cap C'$ and their compact neighborhoods $N_{\lambda'_u}\cap C', g_{_{C,C'}}(N_{\lambda_u}\cap C)$ are subpolyhedra of the annulus  $C'$. From~[\onlinecite[Theorem 4.11, Ch.4]{RuSa}] it follows that there exists a homeomorphism $h_{_{C', C'}}: C'\to C'$ such that $h_{_{C',C'}}(g_{_{C,C'}}(N_{\lambda_u}\cap C))=N_{\lambda'_u}\cap C', h_{_{C',C'}}(\lambda'_u\cap C')=\lambda'_u\cap C'$. It follows from Lemma~\ref{dehn} that there exists a homeomorphism $\tilde h_{_{C', C'}}: C'\to C'$ such that $\tilde h_{_{C',C'}}(g_{_{C,C'}}(N_{\lambda_u}\cap C))=N_{\lambda'_u}\cap C', \tilde h_{_{C',C'}}(\lambda'_u\cap C')=\lambda'_u\cap C'$ and $\tilde h_{_{C', C'}}(g_{_{C,C'}}(\tilde \lambda_u\cap C))=\tilde \lambda'_u\cap C$. Then the superposition $\eta_{_{C,C'}}=\tilde h_{_{C', C'}}g_{_{C', C'}}$ is the desired homeomorphism.

{\it Step 2. Extend the homeomorphism $\eta_{_{C,C'}}$ to the set $N_{\Lambda_s}$}.

Recall that from the definition of neighborhoods $N_{\Lambda_s}, N_{\lambda_u}$ it follows that there exists a homeomorphism $$\xi: \mathbb S^2\times [-1,1]\to N_{\Lambda_s}$$ such that $\xi(\mathbb S^2\times\{0\})=\Lambda_s$, and for each connected component $X$ of the set $N_{\Lambda_s}\cap N_{\lambda_u}$ there is a ball $B^2\subset\mathbb S^2$ such that $X=\xi(B^2\times [-1,1])$, while $\xi(O_i\times [-1,1])\subset \lambda_u$ for some interior point $O$ of the ball $B^2$. Moreover, without loss of generality, we can assume that $\xi(K\times [-1,1])\subset \widetilde \lambda_u$ for some point $K\in \partial B^2$. Denote by $\xi': \mathbb S^2\times [-1,1]\to N_{\Lambda'_s}$ a homeomorphism with similar properties.

We define the homeomorphism $$\theta: \mathbb S^2\times \{-1,1\}\to \mathbb S^2\times \{-1,1\}$$ by $$\theta={\xi'}^{-1}\eta_{_{C,C'}}\xi|_{\mathbb S^2\times \{-1,1\}}.$$ Since each of the mappings $\xi, \xi', \eta_{_{C,C'}}$ is a homeomorphism of the annulus, the homeomorphism $\theta$ either simultaneously preserves or simultaneously reverses the orientation of both spheres $\mathbb S^2\times \{-1\}, \mathbb S^2\times \{1\}$. It follows that there exists an isotopy $\theta_t:\mathbb S^2\to \mathbb S^2$ such that $\theta_0=\theta|_{\mathbb S^2\times \{-1\}}, \theta_1=\mathbb S^2\times \{1\}$. We define the homeomorphism $\Theta: \mathbb S^2\times [-1,1]\to \mathbb S^2\times [-1,1]$ by setting for $x\in \mathbb S^2, y\in [-1,1]$

$$\Theta(x,y)=(\theta_{(y+1)/2}(x),y).$$

Now the desired homeomorphism $\eta: \Sigma\to \Sigma'$ is defined by the formula

\begin{equation}
\eta(p)=\begin{cases}\eta_{_{C,C'}}(p),\ p\in C;\cr
\xi'(\Theta(\xi^{-1}(p))),\ p\in N_{\Lambda_s}.\end{cases}
\end{equation}

\end{proof}





\subsection{Construction of consistent canonical neighborhoods of saddle equilibria}

In this section, we prove a Proposition~\ref{Step1} establishing that the handle decomposition of  the ambient manifold  determinated by the energy  function  of the flow $f^t\in G(1,2,k)$ can be adjusted so that the foliation determinated  by the structure of the direct product  $\mathbb B^i\times \mathbb B^{4-i}$ on each handle  became  $f^t$-invariant as it is defined by item 3 of the following definition.

\begin{figure}
\centering{\includegraphics[width=0.47\textwidth]{canon2.eps}}
\caption{a) canonical neighborhood of a saddle; b) consistent canonical neighborhoods of saddles $\sigma_1,\sigma_2$}
\label{agr}
\end{figure}

\begin{definition}\label{can} A compact neighborhood $V_i$ of a saddle equilibrium state $\sigma_i$ is called a {\it canonical compact neighborhood} if
there exist continuous maps $\pi_{\sigma_i,u}: V_{i}\to W^u_{\sigma_i}, \pi_{\sigma_i,s}: V_{i}\to W^s_{\sigma_i}$ with the following properties (see Fig.~\ref{agr}, a)).

\begin{enumerate}
\item for any points $x\in V_{i}\cap W^u_{\sigma_i}, y\in V_{i}\cap W^s_{\sigma_i}$ the sets $B^s_{\sigma_i,x}=\pi_{\sigma_i,u}^{-1}(x), B^u_{\sigma_i,y}=\pi_{\sigma_i,s}^{-1}(y)$ are smoothly embedded compact balls of dimension $4-i, i$, respectively that intersect each other  transversely;
\item $B^s_{\sigma_i,\sigma_i}\subset W^s_{\sigma_i}$, $B^u_{\sigma_i,\sigma_i}\subset W^u_{\sigma_i}$;
\item $f^{t}(B^s_{\sigma_i,x})\subset B^s_{\sigma_i,f^{t}(x)}$, $f^{t}(B^u_{\sigma_i,y})\supset B^u_{\sigma_i,f^{t}(y)}$.
\end{enumerate}
\end{definition}

The mappings $\pi_{\sigma_i,u}, \pi_{\sigma_i,s}$ define a direct product structure $B^{u}_{\sigma_i, \sigma_i}\times B^s_{\sigma_i, \sigma_i}$ in a neighborhood of $V_{i}$, that is, each neighborhood of $V_i$ is an $i$-handle. As before, we will denote by $P_i$ the foot of the handle $V_i$ and by $Q_i$ the closure of the complement to  $P_i$ in $\partial V_i$.

\begin{definition} Let $f^t\in G_{1,2,k}$. Canonical neighborhoods $V_1, V_2$ of saddle equilibria $\sigma_1, \sigma_2\in \Omega_{f^t}$ are called consistent if either $V_1\cap V_2=\emptyset$ or the following inclusions hold for any point $p\in V_1\cap V_2$ (see Fig.~\ref{agr},b))
$$B^u_{\sigma_1, \pi_{\sigma_1,s}(p)}\subset B^u_{\sigma_2, \pi_{\sigma_2,s}(p)}, B^s_{\sigma_2, \pi_{\sigma_2,u}(p)}\subset B^s_{\sigma_1, \pi_{\sigma_1,u}(p)}.$$
\end{definition}

\begin{propos}\label{Step1} For any flow $f^t\in G_{1,2,k}$, there exists a handle decomposition  of the ambient  manifold $M^4$ such that each handle contains a unique equilibrium state of  $f^t$ and the index of any  handle  is equal to the Morse index of the corresponding equilibrium state. Moreover,
handles of indices 1 and 2 are consistent canonical neighborhoods of saddle equilibria.
\end{propos}
\begin{proof}
We prove the proposition for the case $k=4$ (the proof for $k=5$ is similar).
From Proposition~\ref{hd} it follows that the manifold $M^4$ admits a handle decomposition $H_0\cup H_1\cup H_2\cup H_4$ such that each handle $H_i$ is a compact neighborhood of an equilibrium state of the flow $f^t$ of the corresponding index. By definition, $\Sigma=\partial (H_0\cup H_1)$.

Denote by $e_i: \mathbb B^i\times \mathbb B^{n-i}\to H_i$ the diffeomorphism defining the handle structure. For a point $p\in H_i$, denote by $\widetilde B^s_{i,p} ( \widetilde B^u_{i,p})$ the leaf of the foliation $e_i(\{x\}\times \mathbb B^{4-i})$  ($e_i(\mathbb B^{i}\times \{y\}$, respectively) passing through $p$. {Note that $\widetilde B^u_{i,\sigma_i}=W^u_{\sigma_i}\cap H_i, \widetilde B^s_{i,\sigma_i}=W^s_{\sigma_i}\cap H_i$}.

Let $\widetilde P_2, \widetilde Q_1$ denote the tubular neighborhoods of the knot $\lambda_u$ and the sphere $\Lambda_s$ in $\Sigma$ that belong to the boundary  of the handles $H_2,H_1$, respectively.
We denote by 
 $\widetilde B^s_{1,p_+},\widetilde B^s_{1,p_-}\in \partial H_0$  smoothly embedded in $\partial H_0$  3-balls that are  connected components of  $H_0\cap H_1$, and set $p_+=\widetilde B^s_+\cap W^u_{\sigma_1}, p_-=\widetilde B^s_-\cap W^u_{\sigma_1}$. Let  $\varepsilon_{0}>0$ be a number such that for any smooth $\varepsilon_0$-close to $\widetilde B^s_{1,\sigma_1}$ 3-ball $B^s_p$ (in the $C^1$ topology) the intersections $B^s_p\cap \widetilde B^u_{1,\sigma_1}$, $B^s_p\cap \Sigma$ are transversal. It follows from the classical $\lambda$-lemma (see~[\onlinecite{PaMe}, Lemma 7.1 and the remark in $\S 7$ of Chapter 2]) that there exists $T(\varepsilon_1)>0$ such that for all $t>T(\varepsilon_0)$ the connected component $B^s_{1,f^{-t}(p_{\pm})}$ of the set $f^{-t}(\widetilde B^s_{1,p_\pm})$ containing the point $f^{-t}(p_\pm)$ is a 3-ball $\varepsilon_0$-close to $\widetilde B^s_{1,\sigma_1}$. Then the union $V_{1,\varepsilon_0}=\bigcup\limits_{t\geq T(\varepsilon_0)+1} (\widetilde B^s_{1,f^{-t}(p_\pm)})\cup \widetilde B^s_{1,\sigma_1})$ forms a compact neighborhood of the point $\sigma_1$, foliated by smoothly embedded 3-balls transversal to $\widetilde B^u_{1,\sigma_1}$ and this foliation satisfies item 3 of the Definition~\ref{can}. Without loss of generality, assume that $\partial \widetilde B^s_{1,p_\pm}\subset \Sigma$ (if this is not the case, then replace each ball $\widetilde B^s_{1,p_\pm}$ with a ball that is a connected component of the intersection $\widetilde B^s_{1,p_\pm}\cap H_1$ containing the point $f^{-t}(p_\pm)$).
Set $Q_1=\partial V_{1,\varepsilon_0}\cap \Sigma$. The boundaries of the balls $B^s_{1,f^{-t}(p_{\pm})}$ form a two-dimensional foliation of the set $Q_1$ transverse to the knot $\lambda_u$.

Now let $\varepsilon_{4}>0$ such that for any smooth $\varepsilon_4$-close to $\widetilde B^u_{2,\sigma_2}$ 2-ball $B^u_p$ the intersection $B^u_p\cap \widetilde B^s_{2,\sigma_2}$ is transversal.
By reducing $\varepsilon_0$, if necessary, we can ensure that the intersection $B^u_p$ with each fiber $B^s_{1,f^{-t}(p)}\subset V_{\varepsilon_0}\cap \Sigma$ of the foliation constructed in the previous paragraph is also transversal.
Applying the $\lambda$-lemma again, we find $T(\varepsilon_4)>0$ such that for any point $z\in W^s_{\sigma_2}\cap \partial H_4$ and for all $t>T(\varepsilon_4)$, the connected component $B^u_{2,f^t(z)}$ of the set $f^{t}(\widetilde B^u_{2,z})$ containing the point $f^{t}(z)$, is a 2-ball $\varepsilon_4$-close to $\widetilde B^u_{2,\sigma_2}$.
Then the collection $V_{2,\varepsilon_4}=\bigcup\limits_{t\geq T(\varepsilon_4)+1,z\in W^s_{\sigma_2}\cap \partial H_4} (B^u_{1,f^t(z)})\cup \widetilde B^u_{2,\sigma_2})$ forms a compact neighborhood of $\sigma_2$, fibered by smoothly embedded 2-balls transversal to $\widetilde B^u_{2,\sigma_2}$ and to each fiber $\{\partial B^s_{1,f^{-t}(p)}\}$ of the foliation constructed in the previous step, and this foliation satisfies item 3 of Definition~\ref{can}.

The boundaries of the 2-balls $B^u_{2,f^t(z)}$ define a one-dimensional foliation in the intersection $V_{2,\varepsilon_4}\cap Q_1$ that is transversal to the two-dimensional foliation of the annulus $Q_1$ constructed in the previous paragraph. The connected components of the set $Q_1\cap V_{2,\varepsilon_4}$ are tubular neighborhoods of compact smooth arcs $Q_1\cap \lambda_u$ regularly embedded in $Q_1$. Since the tubular neighborhood is unique, it follows that the one-dimensional  foliation defined on $Q_1\cap V_{2,\varepsilon_4}$ extends to a one-dimensional foliation on the whole of $Q_1$ that is transversal to the previously defined two-dimensional foliation. Let $D^u_x$ denote the fiber of this one-dimensional foliation passing through the point $x\in \Lambda_s$.

Let $\varepsilon_{1}>0$ be such that for any smooth 1-ball $B^u_p$ close to $\widetilde B^u_{1,\sigma_1}$ the intersection $B^u_p\cap \widetilde B^s_{1,\sigma_1}$ is transverse. By reducing $\varepsilon_0$, if necessary, we can ensure that the intersection of $B^u_p$ with each fiber $B^s_{1,f^{-t}(p)}\subset V_{\varepsilon_0}$ of the two-dimensional foliation constructed above is also transverse. Applying the $\lambda$-lemma, we find $T(\varepsilon_1)>0$ such that for any point $x\in \Lambda_s$ and for all $t>T(\varepsilon_1)$, the connected component $B^u_{1,f^t(x)}$ of the set $f^{t}(D^u_{x})$ containing the point $f^{t}(x)$ is a 1-ball $\varepsilon_1$-close to $\widetilde B^u_{1,\sigma_1}$. Then the collection $V_{\varepsilon_1}=\bigcup\limits_{t\geq T(\varepsilon_1)+1} (B^u_{1,f^t(x)})\cup \widetilde B^u_{1,\sigma_1})$ forms a compact neighborhood of $\sigma_1$, foliated by smoothly embedded 1-balls transversal to $\widetilde B^s_{1,\sigma_1}$ and to $\{B^s_{1,f^{-t}(x)}\}$, and this foliation satisfies the Definition~\ref{can}. Set $V=V_{\varepsilon_0}\cap V_{\varepsilon_1}$, choose $T>0$ such that $f^T(V)\cap H_0\subset (\widetilde B^s_+\cup \widetilde B^s_-)$, $f^T(V)\cap \Sigma\subset Q_1$, $f^T(V)\subset {\rm int}\,(H_0\cup H_1)$,
and set $V_{\sigma_1}=f^T(V)\cap H_1$. Since any iteration of a canonical neighborhood is also a canonical neighborhood, and the intersection of $V_{\sigma_1}$ with $\partial H_0, \Sigma$ consists of leaves belonging to the constructed $f^t$-invariant foliations, $V_{\sigma_1}$ is a canonical neighborhood of the saddle $\sigma_1$.

Similarly, we define an $f^t$-invariant foliation $\{B^s_{2,\sigma_2}\}$ of a neighborhood $V_{2,\varepsilon_4}$ of the saddle equilibrium state $\sigma_2$ that is transverse to the foliation $\{B^u_{2,\sigma_2}\}$, and define a canonical neighborhood $V_{\sigma_2}$ that is consistent with $V_{\sigma_1}$ and has a non-empty intersection with $\Sigma$ and $\partial H_0$. The union $N=H_0\cup V_{\sigma_1}\cup V_{\sigma_2}$ is a manifold with boundary $\partial N$ lying in $W^u_\alpha\setminus \alpha$ and being a global secant for the flow $f^t|_{W^u_{\alpha}\setminus \alpha}$. It follows that $\partial N$ is a three-dimensional sphere bounding the ball $B^4_\alpha\subset W^u_\alpha$ (see, for example,~[\onlinecite[Proposition~10]{GuSa25}]). We choose this ball as a handle of index $4$. Then the union $H_0\cup V_{\sigma_1}\cup V_{\sigma_2}\cup B^4_\alpha$ is the desired handle decomposition of the manifold $M^4$.
\end{proof}

\subsection{Local topological conjugacy of flows $f^t, {f'}^t\in G_{1,2,k}$ in consistent canonical neighborhoods of saddle equilibria}

In this section, Proposition~\ref{local_conj} constructs a homeomorphism that conjugates flows $f^t, {f'}^t$ in the union of consistent canonical neighborhoods of their saddle equilibria.

Let $\widehat \Sigma$ be the boundary of the union of handles of indices 0,1 in the handle decomposition described in Proposition~\ref{Step1}, $\widehat \Lambda_s=\widehat \Sigma\cap W^s_{\sigma_1}, \widehat \lambda_u=\widehat \Sigma\cap W^u_{\sigma_2}$. Then $\widehat \Sigma$ is a piecewise smooth global cross-section (in the sense of definition~\ref{seq}) for the restriction of $f^t$ to the set $M^4\setminus ({\rm cl}\, W^u_{\sigma_1}\cup {\rm cl}\, W^s_{\sigma_2})$. It follows from Proposition~\ref{indep} that the set $\{\widehat \Sigma, \widehat \Lambda_s, \widehat \lambda_u\}$ is a scheme of $f^t$ equivalent to the scheme $\{\Sigma, \Lambda_s, \lambda_u\}$. In the following constructions of this section, we will use the set $\{\widehat \Sigma, \widehat \Lambda_s, \widehat \lambda_u\}$ as the flow scheme, omitting the symbol $\ \widehat{}\ $ in the notations.

The consistent canonical neighborhoods, projections, and other objects for the saddle equilibrium state $\sigma'_i$ of the flow ${f'}^t$, similar to those constructed above for the flow $f^t$, will be denoted by $V'_{i}$, $\pi_{\sigma'_i,u}: V'_{i}\to W^u_{\sigma'_i}, \pi_{\sigma'_i,s}: V'_{i}\to W^s_{\sigma'_i}$, etc.

Note that the components $P_2, Q_2$ of the boundary of the handle $V_2$ are solid tori. Recall that the {\it meridian} of a solid torus $P$ is a simple closed curve $\mu\in \partial P$ that is non-homotopic to zero in $\partial P$ and bounds a disk in $P$.

\begin{propos}\label{local_conj}
Let $\{V_1,V_2\}, \{V'_1, V'_2\}$ be pairs of consistent canonical neighborhoods of saddle equilibria of flows $f^t, {f'}^t\in G_{1,2,k}$ whose schemes are equivalent via the homeomorphism $h:\Sigma\to \Sigma'$. Then there exist homeomorphisms $\chi_{1}: V_{1}\to V'_{1}, \chi_{2}: V_{2}\to V'_{2}$ with the following properties:

\begin{enumerate}
\item $\chi_{i}f^t={f'}^t\chi_{i}$ for all $t$ for which the right-hand and left-hand sides of the equality are defined;
\item $\chi_2(p)=\chi_1(p)$ for all points $p\in V_1\cap V_2$;
%\item $\chi_{i}(\lambda_u\cap \Lambda_s)=\lambda'_u\cap \Lambda'_s$;
\item $\chi_{2}$ maps the meridians of the solid tori $P_2, Q_2$ to the meridians of the solid tori $P'_2$, $Q'_2$, respectively.
%\item commutation with projections ??
\end{enumerate}
\end{propos}
\begin{proof}

We describe the construction of the homeomorphisms $\chi_1, \chi_2$ for the case when $V_1\cap V_2\neq \emptyset$ (in the case $V_1\cap V_2=\emptyset$, the constructions are similar). From the equivalence of the schemes  it follows that there exists a homeomorphism $h:\Sigma\to \Sigma$ such that $h(\Lambda_s)=\Lambda'_s, h(\lambda_u)=\lambda'_u$.

Recall that, by the definition of consistent neighborhoods,
$\lambda_u=\partial B^u_{\sigma_2,\sigma_2},\Lambda_s=\partial B^s_{\sigma_1, \sigma_1}$, and $P_2=\Sigma\cap V_{\sigma_2}, Q_1=\Sigma\cap V_{\sigma_1}$ are tubular neighborhoods of the knote $\lambda_u$ and the sphere $\Lambda_s$ in $\Sigma$, respectively. Similar relations are true for the consistent neighborhoods $V'_1, V'_2$ of the saddle equilibrium states of the flow ${f'}^t$.

Let $g^u_1:\partial B^u_{\sigma_1,\sigma_1}\to \partial {B'}^u_{\sigma'_1, \sigma'_1}$, $g^s_2:\partial B^s_{\sigma_2,\sigma_2}\to \partial {B'}^s_{\sigma'_2, \sigma'_2}$ be  arbitrary homeomorphisms.

We associate a time $t_{x} (t_y)\in \mathbb R$ to any  point $x\in B^u_{\sigma_1, \sigma_1}$ $(y\in B^s_{\sigma_2,\sigma_2})$  such that $f^{t_{x}}(x)\in \partial B^u_{\sigma_1,\sigma_1} (f^{t_{y}}(y)\in \partial B^s_{\sigma_2,\sigma_2}$) and define homeomorphisms  

$$h^u_1: B^u_{\sigma_1,\sigma_1}\to B^u_{\sigma'_1,\sigma'_1},\ h^s_2: B^s_{\sigma_2,\sigma_2}\to B^s_{\sigma'_2,\sigma'_2}$$ by


$$h^u_1(x)=\begin{cases}
{f'}^{-t_{x}}(g^u_1(f^{t_{x}}(x))), x\in B^u_{\sigma_1,\sigma_1}\setminus \sigma_1 \cr
\sigma'_1, x=\sigma_1;
\end{cases}$$

$$h^s_2(y)=\begin{cases}
{f'}^{-t_{y}}(g^s_2(f^{t_{y}}(y))), y\in B^s_{\sigma_2,\sigma_2}\setminus \sigma_2\cr
\sigma'_2, x=\sigma_2;
\end{cases}$$

Let $p\in \lambda_u\cap \Lambda_s$. Then $B^u_{\sigma_1, p}\subset \lambda_u=\partial B^u_{\sigma_2,\sigma_2}, B^s_{\sigma_2, p}\subset \Lambda_s=\partial B^s_{\sigma_1,\sigma_1}$. For each point $p\in \lambda_u\cap \Lambda_s$ we define homeomorphisms

$$\xi^u_p: B^u_{\sigma_1, p}\to B^u_{\sigma'_1, h(p)}, \xi^s_p: B^s_{\sigma_2, p}\to B^s_{\sigma'_2, h(p)}$$ such that

$$h^u_1\pi_{u,\sigma_1}=\pi_{u,\sigma'_1}\xi^u_p,\ h^s_2\pi_{s,\sigma_2}=\pi_{s,\sigma'_2}\xi^s_p.$$


From the equivalence of the schemes it follows that the homeomorphisms $\xi^u_p, \xi^s_p$ extend to homeomorphisms


$$g^u_2:\partial B^u_{\sigma_2,\sigma_2}\to \partial {B'}^u_{\sigma'_2, \sigma'_2}, g^s_1:\partial B^s_{\sigma_1,\sigma_1}\to \partial {B'}^s_{\sigma'_1, \sigma'_1},$$ such that for any point $p\in \lambda_u\cap \Lambda_s$

$$g^u_2|_{B^u_{\sigma_1, p}}=\xi^u_p|_{B^u_{\sigma_1, p}},\
g^s_1|_{B^s_{\sigma_2, p}}=\xi^s_p|_{B^s_{\sigma_2, p}}.$$

Let us extend homeomorphisms $g^u_2, g^s_1$ to homeomorphisms
$$h^u_2: B^u_{\sigma_2,\sigma_2}\to {B'}^u_{\sigma'_2, \sigma'_2}, h^s_1: B^s_{\sigma_1,\sigma_1}\to {B'}^s_{\sigma'_1, \sigma'_1}$$

similar to the construction of homeomorphisms $h^u_1, h^s_2$.

Finally, we define the desired homeomorphisms $\chi_{i}: V_i\to V'_i, i\in \{1,2\},$ by the relation
$$\chi_{i}(z)=\pi_{\sigma'_i,u}^{-1}(h^u(\pi_{\sigma_i,u}(z)))\cap {\pi_{\sigma'_i,s}}^{-1}(h^s(\pi_{\sigma_i,s}(z))),\ z\in V_{i}.$$

Note that the solid tori $P_2, Q_2\subset \partial V_2$ have the structure of direct products $\partial B^u_{\sigma_2, \sigma_2}\times B^s_{\sigma_2, \sigma_2}$ and $B^u_{\sigma_2, \sigma_2}\times \partial B^s_{\sigma_2, \sigma_2}$. Then for an arbitrary pair of points $x\in \partial B^u_{\sigma_2,\sigma_2}$, $y\in \partial B^s_{\sigma_2, \sigma_2}$ the curves $\mu_x=\{x\}\times \partial B^s_{\sigma_2, \sigma_2}$, $\mu_y=\partial B^u_{\sigma_2, \sigma_2}\times \{y\}$ are the meridians of $P_2, Q_2$, respectively. It follows from the construction that the curves $\chi_2(\mu_x)=\mu_{h^u_2(x)}, \chi_2(\mu_y)=\mu_{h^s_2(y)}$ are the meridians of the solid tori $P'_2, Q'_2$, respectively.
\end{proof}

\subsection{End of the proof of the Theorem~\ref{cond}}

\begin{propos}\label{Step4} Let  schemes of flows $f^t, {f'}^t\in G_{1,2,k}$ be equivalent. Then there exists a homeomorphism $H: M^4\to M^4$ such that $Hf^t={f'}^tH$.
\end{propos} We prove the proposition for $k=4$ (the arguments for $k=5$ are similar). From Proposition~\ref{Step1} it follows that the sphere $S^4$ admits a decomposition into handles in which the handles of indices 1 and 2 are compact canonical neighborhoods $V_1, V_2 (V'_1, V'_2)$ of saddle equilibrium states of the flow $f^t ({f'}^t$, respectively). Recall that $\partial V_i=P_i\cup Q_i$, where $P_i$ is the foot of the handle $V_i$. Then $Q_1, P_2$ are tubular neighborhoods of the sphere $\Lambda_s$ and the knot $\lambda_u$ in the characteristic cross-section $\Sigma$.

Denote by $\widetilde \lambda_u$ a simple closed curve belonging to $\partial P_2\cup \partial Q_2$, which is a meridian of the solid torus $Q_2$ (and, therefore, a parallel to the solid torus $P_2$). From Proposition~\ref{Step0} it follows that there exists a homeomorphism $\eta:\Sigma\to \Sigma'$ such that $\eta(Q_1)=Q'_1, \eta(P_2)=P_2'$, $\eta(\Lambda_s)=\Lambda'_s, \eta(\lambda_u)=\lambda'_u$, and $\eta(\widetilde \lambda_u)=\widetilde \lambda'_u$.
We define the homeomorphism $\psi: P'_2\to P'_2$ by the formula $\psi=\chi_2\eta^{-1}|_{P'_2}$. The homeomorphism $\psi$ is a homeomorphism of the solid torus onto itself, and therefore preserves the meridian class. Moreover, $\psi(\widetilde \lambda'_u)=\widetilde \lambda'_u$. This implies that the restriction $\psi|_{\partial P'_2}$ is isotopic to the identity map.

Let $N'$ be the collar of $\partial P'_2$ in $\Sigma\setminus {\rm int}\,P'_2$ such that the fibers of the one-dimensional foliation defining the structure $\partial P'_2\times [0,1]$ on $N'$ lie on the two-dimensional spheres foliating the neighborhood $Q'_1$ of $\Lambda'_s$ (the existence of such a collar is guaranteed by the transversality of the intersection of $\Lambda'_s$ with $\lambda'_u$ and~\cite[Ch.4, Section 6, Exercise~2]{Hi}). Using $N'$ and an isotopy to the identity map, we extend $\psi$ to a map $\Psi:\Sigma'\to \Sigma'$ that is the identity outside $N'\cup Q'_2$ and preserves $Q'_1$. Finally, we define a homeomorphism $\theta: \Sigma\to \Sigma'$ by $\theta=\Psi\eta$.

To each point $z\in S^4\setminus \Omega_{f^t}$ we assign a time $t_z\in \mathbb R$ such that $f^{t_z}(z)\subset \Sigma$ and set

\begin{equation} G(z)=\begin{cases}
{f'}^{-t_z}(\theta(f^{t_z}(z))), z\in \Sigma_{f^t};\cr
{f'}^{-t_z}(h_{\sigma_i,\sigma'_i}(f^{t_z}(z))), z\in V_{\sigma_i}, i\in \{1,2\}.
\end{cases} \label{bij}
\end{equation}

Formula~(\ref{bij}) defines a continuous bijection on the set $S^4\setminus \Omega_{f^t}$, which has a unique continuous extension to $\Omega_{f^t}$ to the desired homeomorphism.

\section{Realization of Topological Equivalence Classes}

Let $\Sigma_k$ denote the direct product of $\mathbb S^2\times \mathbb S^1$ for $k=4$ and the sphere $\mathbb S^3$ for $k=5$; $\Lambda, \lambda\subset \Sigma_k$ be a smoothly embedded two-dimensional sphere and a trivial knot, respectively. We call  the set $\{\Sigma_k, \Lambda, \lambda\}$  an {\it abstract scheme} if $\Lambda, \lambda$ intersect each other  transversely and in the case $k=4$ the sphere $\Lambda$ does not divide $\Sigma_4$.

\begin{lemma}\label{rea} For any abstract scheme $\{\Sigma_k, \Lambda, \lambda\}$
there exists a flow $f^t\in G_{1,2,k}, k\in \{4,5\}$, whose scheme is equivalent to the abstract scheme.  
\end{lemma}
\begin{proof} We define a vector field on the handle $H^4_i=\mathbb B^i\times \mathbb B^{4-i}$   by the system $$\begin{cases}\dot{x}=x,\cr \dot y =-y,\end{cases}$$ where $x\in \mathbb B^i, y\in \mathbb B^{4-i}$, $i\in \{0,\dots, 4\}$. In case $k=4(5)$, we glue the handle $H^4_1$ to one (two, respectively) copies of the handle $H^4_0$ using a diffeomorphism that ensures smooth gluing of the vector fields. By smoothing the resulting manifold in the standard way, we obtain a manifold $Y$ whose boundary is diffeomorphic to $\Sigma_k$, and a vector field on it transversal to $\partial Y$ which has two (three) singular points: one (two, respectively) sinks and a saddle of index 1. The stable manifold of the saddle intersects $\partial Y$ along the two-dimensional sphere $\Lambda_s$. In the case $k=4$ this sphere does not divide $\partial \Sigma$. Then there exists a diffeomorphism $h: \partial Y\to \Sigma_k$ such that $h(\Lambda_s)=\Lambda$. We set $\lambda_u=h^{-1}(\lambda)$. We attach the handle $H^4_2$ to $Y$ by means of  a diffeomorphism that maps the foot sphere of the handle  to the knot $\lambda_u$. Since the knot $\lambda_u$ is trivial, the boundary of the resulting manifold $Z$ is diffeomorphic to the sphere $\mathbb S^3$. We attach the handle $H^4_4$ to $Z$, resulting in a closed smooth manifold $M^4$ and a vector field on it which  induces a flow $f^t\in P$ whose scheme is equivalent to the given abstract scheme.
\end{proof}

\subsection{Proof of Theorem~\ref{realiz}}

The first part of Theorem~\ref{realiz} follows immediately from Lemma~\ref{rea}. We show that for any integer $\gamma \geq 1$, the class $G_{1,2,4}$ contains a countable set of topologically nonequivalent flows with exactly $(2\gamma+1)$ heteroclinic trajectories. To do this, by Lemma~\ref{rea} and Theorem~\ref{cond}, it suffices to provide a countable set of pairwise nonequivalent abstract schemes in each of which the intersection of the knot $\lambda$ and the sphere $\Lambda$ consists of exactly $(2\gamma+1)$ points. Examples of such abstract schemes are shown in Fig.~\ref{diff}.

\begin{figure}
\centering{\includegraphics[width=0.4\textwidth]{het33.eps}}
\caption{Nontrivial sphere $\Lambda$ and trivial knot $\lambda_{k,i}$ in $\Sigma=\mathbb S^2\times \mathbb S^1$}
\label{diff}
\end{figure}

A  top part of Figure~\ref{diff} shows three cases of mutual arrangement of a homotopically nontrivial two-dimensional sphere $\Lambda$ and trivial knot $\lambda_{3,0}, \lambda_{3,1}, \lambda_{3,2}$ in $\Sigma_4=\mathbb S^2\times \mathbb S^1$, transversely intersecting at three points. To obtain $\mathbb S^2\times \mathbb S^1$, it is necessary to glue together   the boundary spheres of the annulus  highlighted in gray  by a homothety. The series of knots shown in the figure naturally extends to a countable series of trivial knots $\{\lambda_{3, i}\}$. We  show that there is no  homeomorphism $h: \Sigma_4\to \Sigma_4$ such that $h(\Lambda)=\Lambda$, $h(\lambda_{3,i})=\lambda_{3,j}$ for $i\neq j$. Assume the contrary.
\begin{figure}
\centering{\includegraphics[width=0.34\textwidth]{het33+.eps}}
\caption{Ball $B=\Sigma\setminus ({\rm int}\, N_{_{\Lambda}}\cup {\rm int}\, C_b)$ and arcs $l_{i,j}, e_{i,j}\subset B$ for $i\in \{0,1\}, j\in \{1,2\}$}
\label{diff+}
\end{figure}

Let $N_{_\Lambda}, N_{\lambda_{3,i}}$ be compact tubular neighborhoods of the sphere $\Lambda$ and the knot $\lambda_{3,i}$ in $\Sigma_4$, respectively. Since the sphere $\Lambda$ does not divide $\Sigma_4$, there exists a compact arc $b_{3,i}\subset \lambda_{3,i}$ whose endpoints lie on the  different connected components of the boundary of the annulus $N_{_\Lambda}$. Let $C_b\subset \lambda_{3,i}$ be a connected component of the set $N_{\lambda_{3,i}}\setminus {\rm int}\,\Lambda$ containing the arc $b_{3,i}$. The set $B=\Sigma_4\setminus {\rm int}\, (N_{_\Lambda}\cup  C_{b})$ is homeomorphic to the ball $\mathbb B^3$. Let $l_{i,1}, l_{i,2}$ denote the connected components of the set $B\cap \lambda_{3,i}$ (see Fig.~\ref{diff+}, which illustrates cases $i\in \{0,1\}$). The endpoints of each of the arcs $l_{i,1}, l_{i,2}$ belong to the boundary $\partial B$ of the ball $B$. We  connect the endpoints of the arcs $l_{i,j}$ by an arc $e_{i,j}\subset \partial B$ such that $e_{i,1}\cap e_{i,2}=\emptyset$ and set  $k_{i,j}=l_{i,j}\cup e_{i,j}$. We glue the ball $\mathbb B^3$ to the ball $B$, arbitrarily identifying the boundary spheres $\mathbb S^2$ and  $\partial B$. As a result, we obtain the sphere $S^3$ and the link $L_{3,i}=\{k_{i,1}, k_{i,2}\}$ in it. By construction, for each $i$, at least one of the knots, say $k_{i,2}$, is trivial, and there exists a three-dimensional ball $D\subset S^3$ such that $k_{i,2}\subset {\rm int}\, D$, $k_{i,1}\cap D=\emptyset$.
The existence of a homeomorphism $h: \Sigma_4\to \Sigma_4$ implies the existence of a homeomorphism $H: S^3\to S^3$ such that $H(k_{i,1})=k_{j,1}$ for $i\neq j$. But the knot $k_{i,1}$ is the connected sum of the trivial knot and $i$ copies of the trefoil. All such knots are pairwise nonequivalent, hence, we get a contradiction.  

The bottom of the figure~\ref{diff} shows how the series of knots $\lambda_{3,i}$ is modified into a series of knots $\lambda_{2\gamma+1,i}$, each of which has a given number $2\gamma+1$ of intersections with the sphere $\Lambda$. Arguments similar to those presented above prove the existence of a countable set of topologically nonequivalent flows from the class $G_{1,2,4}$ with a given number $2\gamma+1$ of heteroclinic intersections. Thus, Theorem~\ref{realiz} is proved.


\begin{acknowledgments}
The work was supported by the Russian Science Foundation (grant No. 23-71-30008).
\end{acknowledgments}


\nocite{*}
\bibliography{aipsamp}% Produces the bibliography via BibTeX.

\end{document}
%
% ****** End of file aipsamp.tex ******
